% Notes-only addendum to the independently audited D.93 candidate.
% No integration or promotion until full independent review of this argument.
\subsubsection{A proposed all-depth integral Smith calculation for D.93}

The following strengthens only the finite presentation in D.93.
It keeps the entire actual arithmetic reduction of that candidate.
For every \(r\ge1\), it proves the proposed exact formula
\[
 \mathcal S_{s_r}\simeq
 \left(\bigoplus_{j=1}^{r-1}\mathbb Z/5^j\right)
       \oplus\mathbb Z/5^{r+1}.                       \tag{D93S.1}
\]
The first sum is zero when \(r=1\). The proof is an integral
extension calculation; it does not infer Smith factors from a
determinant or from the first two numerical examples.

\paragraph{An exact extension by one order-five factor.}
Put \(n=5^r\), \(R=\mathbb Z_5[t]/(t^n-1)\), \(u=t-1\),
and \(I=uR\). Let \(\sigma_\chi(t)=t^\chi\).
On the same underlying free module \(R\), distinguish
\[
 \mathscr A:\ \chi\sigma_\chi,\qquad
 \mathscr B:\ \chi S_\chi(t)^3\sigma_\chi,
       \qquad S_\chi(t)=\frac{t^\chi-1}{t-1}.
\]
The module \(\mathscr B_\Gamma\), with
\(\Gamma=\mathbb Z_5^\times\), is exactly the compact kernel
in D.93. Write
\[
                  \mathcal N=1+t+\cdots+t^{n-1}.
\]
Multiplication by \(u^3\) gives the exact equivariant sequence
\[
 0\longrightarrow \mathbb Z_5\mathcal N(\chi^4)
   \longrightarrow \mathscr B
   \xrightarrow{u^3} I^3\mathscr A
   \longrightarrow0.                                \tag{D93S.2}
\]
Indeed the equivariance is
\(u^3 S_\chi(t)^3=\sigma_\chi(u^3)\).
The kernel of multiplication by \(u^3\) in \(R\) is precisely
\(\mathbb Z_5\mathcal N\): over a characteristic-zero splitting
field it is supported only at \(t=1\), hence is killed by \(u\);
the coefficient equations for \(u f=0\) then say that all
coefficients of \(f\) are equal. This also proves the integral
kernel assertion. The action on this generator is \(\chi^4\).

The quotient \(\mathscr A/I^3\mathscr A\) has the three
successive tame characters \(\chi,\chi^2,\chi^3\), on its
augmentation filtration. These assertions include any
five-torsion in its two positive-degree pieces. All three
characters are nontrivial on the order-four subgroup of
\(\Gamma\). Exact tame averaging therefore kills every
\(\Gamma\)-homology group of the quotient. Consequently
\[
 (I^3\mathscr A)_\Gamma\simeq\mathscr A_\Gamma,
 \qquad H_1(\Gamma,I^3\mathscr A)=H_1(\Gamma,\mathscr A)=0.
                                                               \tag{D93S.3}
\]
For the last equality, the completed procyclic resolution
after tame averaging identifies \(H_1\) with invariants.
The action of two on \(\mathscr A\) is twice a finite
permutation, so its fixed module is zero. All modules here
are finite free compact \(\mathbb Z_5\)-modules before the
displayed finite quotients; the compact exact resolution is
being used with its actual topology.

Decompose the basis \(t^a\) of \(\mathscr A\) into its orbits
under \(a\mapsto2a\) in \(\mathbb Z/n\).
The zero orbit has weight two and contributes nothing.
For each \(1\le j\le r\), the orbit of
\(t^{5^{r-j}}\) has length \(\ell_j=4\cdot5^{j-1}\);
its weighted-cycle cokernel is
\(\mathbb Z_5/(2^{\ell_j}-1)\simeq\mathbb Z/5^j\).
The exact homology sequence of (D93S.2) now gives
\[
 0\longrightarrow\mathbb Z/5\,[\mathcal N]
 \longrightarrow\mathscr B_\Gamma
 \longrightarrow\bigoplus_{j=1}^r\mathbb Z/5^j
 \longrightarrow0.                                  \tag{D93S.4}
\]
In particular the norm class has exact order five. It is
not discarded by a rational calculation.

\paragraph{An explicit lift of the largest orbit generator.}
Let \(\sigma=\sigma_2\), and write
\[
 A=2\sigma,\qquad B=2(1+t)^3\sigma
\]
for the actual operators on \(\mathscr A,\mathscr B\).
The integral-at-five polynomials
\[
 G=-1-\frac{u}{3}+\frac{2u^2}{21},
       \qquad f=\frac4{21}(t+3)
\]
satisfy the literal polynomial identity
\[
                 u^3 f=t+(A-1)G.                     \tag{D93S.5}
\]
Thus \([f]\) maps in (D93S.4) to the largest orbit
generator \([t]\), up to the stated natural identification.

Put
\[
 \ell=4\cdot5^{r-1},\qquad
 c=2^\ell-1,\qquad
 Q=\sum_{i=0}^{\ell-1}2^i t^{2^i},\qquad H=Q+cG .
                                                               \tag{D93S.6}
\]
Here \(v_5(c)=r\) and \(c/n\equiv3\pmod5\).
The weighted orbit sum gives
\[
                       (A-1)H=c\,u^3f
                          \quad\hbox{in }R.           \tag{D93S.7}
\]
Indeed \(A^\ell t=2^\ell t^{2^\ell}=(1+c)t\) in
\(R\), since \(n\mid c\). This argument is in the actual
finite group ring; it does not set \(t^{2^\ell}=t\)
in the unrestricted polynomial ring.

We give an integral lift of \(H\) through multiplication
by \(u^3\). Write \(h_i\) for its coefficient of \(u^i\)
in the unrestricted polynomial ring. Its constant
coefficient is zero. For
\[
        S_k=\sum_{i=0}^{\ell-1}2^{ki}
               =\frac{(1+c)^k-1}{2^k-1},
\]
its first three coefficients are
\[
 \begin{aligned}
 h_1&=\frac{c+c^2}{3},\\
 h_2&=-\frac c{42}+\frac{c^2}{21}+\frac{c^3}{14},\\
 h_3&=\frac{S_4-3S_3+2S_2}{6}.
 \end{aligned}                                       \tag{D93S.8}
\]
The first two lie in \(n\mathbb Z_5\). Therefore
\[
 W_0=h_1/n,\qquad
 W_1=\bigl(h_2-\tbinom n2W_0\bigr)/n
                       \quad\hbox{belong to }\mathbb Z_5,
\]
and there is an actual polynomial \(Z\in\mathbb Z_5[u]\)
with
\[
 H-(t^n-1)(W_0+W_1u)=u^3Z.                            \tag{D93S.9}
\]
Since \(v_5(\binom n2)=v_5(\binom n3)=r\), its
constant coefficient \(z_0=Z|_{u=0}\) satisfies
\[
                         z_0\equiv h_3\pmod{5^r}.      \tag{D93S.10}
\]

\paragraph{The norm extension is nonzero on the largest factor.}
Equations (D93S.7)--(D93S.9), and the integral kernel
calculation in (D93S.2), give
\[
             c f-(B-1)Z=k\mathcal N
                   \quad\hbox{in }R,\quad k\in\mathbb Z_5.
                                                               \tag{D93S.11}
\]
Evaluation at \(t=1\) yields the exact equality
\[
                 nk=c\,\frac{16}{21}-15z_0.            \tag{D93S.12}
\]
We compute \(k\) modulo five. By (D93S.10), replacing
\(z_0\) by \(h_3\) changes \(15z_0/n\) by a multiple
of five. The terms \(S_2,S_3\) in (D93S.8) have
valuation at least \(r\), so they disappear after this
operation modulo five. The \(S_4\) term gives exactly
\[
 \frac{15h_3}{n}
   \equiv\frac{(1+c)^4-1}{6n}
   \equiv\frac{2c}{3n}\pmod5.
\]
All expressions compared here are integral at five.
Consequently
\[
       k\equiv\frac cn\left(\frac{16}{21}-\frac23\right)
        =\frac{2c}{21n}\equiv1\pmod5.                 \tag{D93S.13}
\]
Thus \(c[f]\) is a nonzero norm class in (D93S.4).
Since \(v_5(c)=r\), the lift \([f]\) has exact order
\(5^{r+1}\). This proves the required nonsplitting over
the largest cyclic quotient.

\paragraph{Splitting every smaller factor.}
Rescale this lift by a five-adic unit to obtain \(x_r\)
with \(5^r x_r=[\mathcal N]\). Choose arbitrary lifts
\(x_j\) of the smaller orbit generators, \(1\le j<r\).
For some \(a_j\in\mathbb Z/5\), they satisfy
\(5^jx_j=a_j[\mathcal N]\).
Choose integral representatives of \(a_j\) and replace
\[
                     x_j\ \longmapsto\
                     x_j-a_j5^{r-j}x_r .
 \]
The new lift has order exactly \(5^j\). On the quotient
in (D93S.4), these replacements are invertible triangular
basis changes: the \(r\)-th generator is kept, and each
smaller one is changed only by a permissible multiple of
that generator. The modified lifts and \(x_r\) generate
the whole module. Their product of orders equals the
order supplied by (D93S.4). They therefore give precisely
the direct-sum decomposition (D93S.1).

This closes the all-depth elementary-divisor calculation.
Its only computational support is the literal symbolic
identities and the already retained small matrix checks;
the extension argument above supplies the proof for
every \(r\).
